% Section 2: Background
% Compactness-VRA Tradeoff Quantification

\section{Background}

\subsection{Compactness in Redistricting Law and Practice}

District compactness---the intuitive notion that legislative districts should be ``reasonably regular'' in shape rather than bizarrely contorted---has been a normative principle in American redistricting since the 19th century. Early state constitutions required districts to be ``compact'' or ``as compact as possible,'' though without defining the term mathematically \citep{young1988}. The practical motivation was twofold: compact districts facilitate constituent-representative communication (minimizing travel distance) and prevent gerrymandering (irregular shapes suggest manipulation for partisan or racial advantage).

The Supreme Court's ``one person, one vote'' doctrine in \emph{Reynolds v. Sims} (1964) established population equality as the primary constitutional requirement, mandating that congressional districts within a state contain equal population (within narrow tolerance). However, equal population alone does not prevent gerrymandering---partisan actors can pack or crack opposition voters into uncompetitive districts while maintaining perfect population balance. Compactness serves as a geometric constraint that limits manipulation: highly compact districts are harder to gerrymander because their shapes are determined by contiguity and area minimization rather than partisan engineering.

Despite widespread normative acceptance, compactness lacks consistent legal enforcement. No federal constitutional requirement mandates compact districts, and the Supreme Court in \emph{Shaw v. Reno} (1993) used compactness primarily as evidence of racial motivation rather than as an independent requirement. State constitutional compactness mandates vary widely: some states require districts to be ``as compact as practicable'' (e.g., Florida, California), others mention compactness alongside competing criteria (e.g., preserving communities of interest, respecting county boundaries), and some provide no compactness guidance at all \citep{levitt2018}.

Quantifying compactness has proven challenging. \citet{polsby1991} popularized the Polsby-Popper score---the ratio of a district's area to the area of a circle with the same perimeter ($PP = 4\pi A / P^2$)---which ranges from 0 (least compact) to 1 (perfect circle). This metric penalizes irregular perimeters regardless of population distribution. Alternative metrics include the Reock score (ratio of district area to smallest enclosing circle), convex hull ratio (ratio to smallest enclosing convex polygon), and graph-theoretic measures based on edge-cut minimization \citep{altman1998}. Courts have generally declined to adopt specific metrics, instead treating compactness as a qualitative judgment informed by visual inspection and expert testimony.

\subsection{Why Compactness Matters: Normative Foundations}

While compactness is widely accepted as a desirable redistricting principle, the \emph{why} is often left implicit. Different normative frameworks justify compactness requirements through distinct mechanisms, and understanding these foundations helps clarify which compactness metrics best operationalize each goal.

\textbf{Travel Distance Minimization}: Compact districts reduce the geographic area representatives must cover, minimizing travel time for constituent services, town halls, and campaign events. This justification dates to pre-automobile eras when representatives literally traveled by horse between communities. While modern transportation has reduced travel costs, minimizing geographic extent remains relevant for resource-constrained campaigns and ensuring representatives can physically reach all constituents. Metrics emphasizing \emph{dispersion} (Reock score, moment of inertia) best capture this concern by measuring how spread out a district is from its center.

\textbf{Community Cohesion}: Compact districts are more likely to encompass communities with shared interests, experiences, and policy priorities. Elongated or irregular districts may artificially combine disparate regions (urban cores with distant suburbs, coastal communities with inland areas) that lack common concerns. This justification treats compactness as a proxy for \emph{communities of interest}---the principle that redistricting should keep together voters who share social, economic, or geographic ties. However, compactness is an imperfect proxy: circular districts may divide cities, while non-compact districts may preserve county boundaries. Metrics emphasizing \emph{contiguity and cohesion} (edge-cut minimization, graph-theoretic measures) align with this goal.

\textbf{Preventing Gerrymandering}: Bizarre district shapes are telltale signs of manipulation for partisan or racial advantage. The infamous ``Gerry-mander'' (1812) earned its name precisely because its salamander-like shape revealed partisan intent. Compactness serves as a \emph{procedural safeguard} against gerrymandering: if districts are required to be compact, mapmakers have less freedom to pack or crack opposition voters. This anti-gerrymandering justification treats compactness as a constraint on discretion rather than a substantive good. Metrics emphasizing \emph{perimeter irregularity} (Polsby-Popper, boundary complexity) best detect manipulation, as gerrymandered districts typically exhibit convoluted boundaries designed to include or exclude specific neighborhoods.

\textbf{Democratic Legitimacy and Transparency}: Compact districts are easier for voters to understand, visualize, and evaluate. When districts follow natural geographic boundaries (rivers, mountain ranges) or administrative units (counties, cities), voters can readily identify ``their'' district and hold representatives accountable. Non-compact districts with complex boundaries obscure accountability and suggest hidden manipulation. This transparency justification treats compactness as enhancing \emph{democratic legitimacy}---voters are more likely to trust redistricting processes that produce simple, understandable districts. Metrics emphasizing \emph{deviation from simple shapes} (convex hull ratio, comparison to regular polygons) capture this transparency concern.

\textbf{Metric Selection and Normative Priorities}: Our analysis uses four compactness metrics (edge cut, Polsby-Popper, Reock, convex hull ratio) precisely because different metrics emphasize different normative goals. Polsby-Popper's perimeter-to-area ratio detects gerrymandering through boundary irregularity. Reock's enclosing circle measures dispersion relevant to travel distance. Edge-cut minimization promotes community cohesion by reducing artificial boundaries. Convex hull ratio assesses deviation from simple, transparent shapes. By reporting results across multiple metrics and demonstrating cross-metric consistency (Table 8: correlations range $r = -0.67$ to $r = +0.82$), we show that our findings are robust to different normative frameworks for why compactness matters.

This multi-metric approach acknowledges that redistricting involves normative trade-offs not just between VRA compliance and compactness, but also between different conceptions of what compactness itself should achieve. Courts and commissions that prioritize anti-gerrymandering may emphasize Polsby-Popper; those prioritizing community cohesion may prefer edge-cut; those prioritizing constituent access may favor Reock. Our analysis demonstrates that edge-weighted optimization generally improves compactness \emph{regardless of which normative framework} one adopts, strengthening the claim that VRA-compactness conflicts are often algorithmic artifacts rather than inherent tensions.

\subsection{The Voting Rights Act and Majority-Minority Districts}

The Voting Rights Act of 1965 transformed redistricting by prohibiting practices that dilute minority voting strength. Section 5 originally required federal preclearance for redistricting changes in covered jurisdictions (primarily Southern states with histories of racial discrimination), while Section 2 (amended in 1982) prohibits vote dilution nationwide. The Supreme Court's \emph{Thornburg v. Gingles} (1986) decision established three preconditions for proving Section 2 violations: (1) the minority group is sufficiently large and geographically compact to constitute a majority in a single-member district, (2) the minority group is politically cohesive, and (3) the white majority votes as a bloc to usually defeat the minority's preferred candidates.

The first \emph{Gingles} precondition---geographic compactness---explicitly links VRA compliance to redistricting geometry. Courts interpret this as requiring that minority populations be ``sufficiently concentrated'' to form districts without ``bizarre'' shapes, but the threshold remains ambiguous. In practice, redistricting plans have created \textbf{majority-minority (MM) districts} where minorities constitute 50\% or more of the population, enabling them to elect their preferred candidates under conditions of racially polarized voting.

However, the creation of MM districts has generated substantial controversy and litigation. \emph{Shaw v. Reno} (1993) established that districts drawn ``so bizarrely on their face that they are unexplainable on grounds other than race'' may violate equal protection, even if intended to comply with the VRA. North Carolina's 12th Congressional District---a narrow 160-mile corridor following Interstate 85 to connect Black voters in Charlotte, Greensboro, Durham, and Winston-Salem---became the canonical example of extreme non-compactness justified by racial representation. Subsequent cases (\emph{Miller v. Johnson} 1995, \emph{Bush v. Vera} 1996) scrutinized MM districts for subordinating traditional redistricting principles (compactness, contiguity, respect for political subdivisions) to racial goals, suggesting that VRA compliance could be achieved only at the cost of other desirable redistricting properties.

This legal history has created a widespread assumption: achieving VRA compliance through MM districts inherently requires sacrificing compactness. Legislatures defend non-compact MM districts as ``necessary'' to satisfy Section 2, courts weigh VRA compliance against compactness on a case-by-case basis, and political opponents of VRA enforcement cite compactness concerns to resist creating MM districts.

\subsection{Defining VRA Section 2 Compliance: Operational Standards}

Section 2 of the Voting Rights Act prohibits electoral practices that ``result in a denial or abridgement of the right... to vote on account of race or color.'' However, the statute does not prescribe specific remedies, district thresholds, or demographic targets. This creates ambiguity: what constitutes ``VRA compliance'' in redistricting? Courts, scholars, and practitioners have adopted varying interpretations, each with different implications for the relationship between minority representation and compactness.

\textbf{Proportional Representation via Majority-Minority Districts}: The most common interpretation treats VRA compliance as requiring districts where minorities constitute $\geq 50\%$ of the population (or voting-age population), with the number of such districts proportional to the state's minority percentage. Under this standard, a state with 40\% Black population and 10 congressional districts should create 4 majority-minority districts. This ``proportional MM'' standard has dominated redistricting litigation in the South, particularly in states covered by former Section 5 preclearance. Our study adopts this interpretation, testing whether edge-weighted optimization can create target numbers of 50%+ MM districts (Alabama: 2/7, Georgia: 5/14, Louisiana: 2/6, Mississippi: 2/4, South Carolina: 3/7) based on proportional representation goals.

\textbf{Alternative Thresholds and District Types}: However, courts have recognized that Section 2 compliance may be satisfied through district types other than 50%+ MM districts:

\begin{itemize}
    \item \textbf{Coalition Districts} (typically 40-45\% minority): Districts where a minority group, though not a numerical majority, can elect preferred candidates through coalitions with sympathetic voters from other groups. \emph{LULAC v. Perry} (2006) recognized that Latino voters in Texas achieved electoral success in districts with 40-45\% Latino population through coalition-building. Coalition districts satisfy *Gingles* prong 1 (``sufficiently large'') without requiring 50%+ thresholds.

    \item \textbf{Influence Districts} (typically 25-40\% minority): Districts where minorities, while unable to elect preferred candidates alone, constitute a large enough share to influence electoral outcomes and policy priorities. \emph{Georgia v. Ashcroft} (2003) acknowledged influence districts as potentially satisfying Section 2 under certain circumstances, though this doctrine has been contested.

    \item \textbf{Crossover Districts} (variable thresholds): Districts where minorities constitute less than 50\% but regularly elect preferred candidates due to substantial white crossover voting. \emph{Bartlett v. Strickland} (2009) held that Section 2 does not \emph{require} creation of crossover districts where minorities are less than 50\%, but left open whether such districts can satisfy Section 2 if created voluntarily.
\end{itemize}

\textbf{Justification for 50\% Threshold in This Study}: We focus on 50%+ majority-minority districts for three reasons: (1) \textbf{Legal precedent}: The 50\% threshold has been the dominant standard in redistricting litigation in our study states (Alabama, Georgia, Louisiana, Mississippi, South Carolina), all of which were covered by Section 5 preclearance and have extensive VRA litigation history involving MM district creation. State-specific consent decrees and legislative practice have established 50%+ as the operative standard in these jurisdictions. (2) \textbf{Conservative approach}: Testing the most stringent threshold (50%+) establishes a lower bound on VRA-compactness tradeoffs. If 50%+ MM districts can be created with minimal or negative compactness costs, then coalition districts (40-45%) are even more feasible, strengthening our findings. (3) \textbf{Gingles precondition 1}: The first *Gingles* requirement---that the minority group be ``sufficiently large and geographically compact to constitute a majority in a single-member district''---explicitly references majority status, suggesting 50%+ as the threshold for this precondition.

\textbf{Sensitivity to Alternative Thresholds}: While our primary analysis uses 50%+ thresholds, our edge-weighting parameter sweeps test minority thresholds ranging from 40\% to 55\%. This allows assessment of how lowering the threshold (to 40-45\%, consistent with coalition districts) affects feasibility and compactness tradeoffs. South Carolina, which fails to achieve 3 MM districts at 50%+ thresholds (feasibility ratio 1.22), may achieve its target through coalition districts at 40-45\% thresholds, with implications for courts applying \emph{Gingles} analysis. Our Results section reports outcomes across this threshold range, enabling readers to assess VRA compliance under alternative interpretations.

\textbf{Population vs. Voting-Eligible Population}: An additional definitional question concerns the population base: should ``majority-minority'' be defined based on total population, voting-age population (VAP), or citizen voting-age population (CVAP)? The Supreme Court's \emph{Evenweil v. Abbott} (2016) held that states \emph{may} use total population for redistricting, but left unresolved whether Section 2 analysis should consider VAP or CVAP. Our primary analysis uses total population (consistent with Census apportionment and *Evenweil*), but Section 6.4 (Limitations) discusses implications of age, citizenship, and registration gaps that may reduce minority \emph{electoral} strength below minority \emph{population} percentages.

By explicitly defining ``VRA compliance'' as achieving target numbers of 50%+ majority-minority districts based on proportional representation, we clarify our operational standard while acknowledging that Section 2 jurisprudence permits alternative interpretations. Our findings---that non-MM districts generally gain compactness when MM districts are created---hold regardless of which threshold one adopts (50%+ MM, 40-45\% coalition, or 25-40\% influence), though the specific feasibility thresholds and Pareto frontier slopes may shift with lower thresholds.

\subsection{Multi-Objective Optimization and Pareto Frontiers}

Redistricting inherently involves multiple competing objectives: equal population (constitutional requirement), compactness (state constitutional/normative preference), respect for communities of interest (cities, counties, neighborhoods), partisan fairness (proportional representation), and minority representation (VRA compliance). Optimizing any single objective is straightforward---maximize compactness by drawing circular districts, maximize partisan advantage by cracking opposition voters, maximize minority representation by concentrating minority voters---but simultaneous optimization of multiple objectives typically produces \emph{tradeoffs} where improving one objective requires degrading another.

\textbf{Pareto efficiency} formalizes this tradeoff structure: a solution is Pareto-optimal if no alternative solution improves one objective without worsening at least one other objective. The set of all Pareto-optimal solutions forms the \textbf{Pareto frontier}, representing the achievable tradeoff space. Plans that lie below the frontier are \emph{dominated}---they are strictly worse than frontier alternatives on at least one objective without being better on others---and are therefore unjustifiable from a multi-objective optimization perspective.

In redistricting, Pareto frontier analysis can reveal whether perceived tradeoffs are inherent (frontier slopes upward, requiring compactness sacrifice for additional MM districts) or artificial (dominated plans exist due to suboptimal algorithms). Several studies have applied Pareto analysis to redistricting: \citet{mehrotra1998} optimized compactness-population tradeoffs using integer programming, \citet{ricca2013} developed multi-objective heuristics for Italian redistricting, and \citet{king2022} explored partisan fairness-compactness tradeoffs using simulated annealing.

However, prior work has not systematically analyzed \textbf{VRA-compactness Pareto frontiers} across multiple states to determine whether the assumed tradeoff is consistent or state-dependent. Existing studies either focus on single states \citep{mccartan2020}, treat VRA compliance as a binary constraint rather than a continuous objective \citep{deford2021}, or ignore VRA requirements entirely \citep{levin2020}. This gap leaves open the question: \emph{Is the VRA-compactness tradeoff an inherent geometric property, or does it vary based on demographic geography?}

\subsection{Graph Partitioning and Demographic-Aware Optimization}

Modern redistricting increasingly relies on graph partitioning algorithms to automate district construction. The redistricting problem can be formulated as partitioning a graph $G = (V, E)$ where nodes $V$ represent census tracts (or blocks) and edges $E$ represent geographic adjacency. The partitioning must satisfy hard constraints (equal population, contiguity) while optimizing soft objectives (compactness, communities of interest).

The METIS algorithm \citep{karypis1998}, originally developed for scientific computing, has become widely used in redistricting due to its speed and effectiveness at minimizing \textbf{edge cut}---the number of edges connecting nodes in different partitions. Minimizing edge cut produces compact districts because it penalizes long, irregular boundaries. METIS uses recursive bisection: repeatedly splitting the graph into two balanced partitions until the desired number of districts is reached.

Standard METIS treats all edges identically, optimizing purely for geometric compactness without regard for demographic characteristics. Recent research has explored incorporating demographic information through \textbf{multi-constraint partitioning}, where nodes have multiple weights (total population, minority population) and partitions must balance both constraints simultaneously \citep{levin2020}. However, multi-constraint approaches rigidly enforce demographic quotas, which may be infeasible or overly restrictive.

An alternative approach---which we employ in this paper---is \textbf{edge-weighting}: assigning higher costs to edges connecting minority-concentrated tracts, making these edges ``expensive'' to cut. This encourages METIS to keep minority communities together while maintaining flexibility in achieving demographic targets. Edge-weighting has been used in other graph partitioning contexts (load balancing, parallel computing) but has not been systematically evaluated for VRA-compliant redistricting.

\subsection{Research Gap: Systematic VRA-Compactness Tradeoff Quantification}

Despite decades of litigation, legislation, and scholarship on redistricting, no prior study has systematically quantified the VRA-compactness tradeoff across multiple states with district-level breakdown. Existing work has either:
\begin{itemize}
    \item \textbf{Qualitatively acknowledged tension} between VRA compliance and compactness without quantitative measurement \citep{pildes1993,issacharoff2002}
    \item \textbf{Studied single states} in isolation, limiting generalizability of findings \citep{mccartan2020}
    \item \textbf{Focused on partisan fairness} rather than racial representation, treating VRA compliance as a secondary constraint \citep{deford2021,king2022}
    \item \textbf{Treated VRA as binary} (compliant vs. non-compliant) rather than as a continuous objective amenable to Pareto frontier analysis \citep{levin2020}
\end{itemize}

This gap leaves critical policy questions unanswered:
\begin{itemize}
    \item \textbf{Are non-MM districts harmed?} Political opposition to VRA enforcement claims that concentrating minorities in MM districts forces irregular shapes on all districts, ``spreading the pain'' to non-minority voters. Is this empirically supported?
    \item \textbf{Can both objectives improve simultaneously?} Do win-win configurations exist where VRA compliance \emph{improves} compactness, or does any VRA optimization necessarily degrade geometric quality?
    \item \textbf{Are tradeoffs universal or state-dependent?} Does the VRA-compactness relationship exhibit consistent slopes across states, or do demographic geography and spatial clustering produce heterogeneous patterns?
    \item \textbf{What determines feasibility thresholds?} At what demographic concentration levels does VRA compliance become geometrically infeasible regardless of algorithmic sophistication?
\end{itemize}

Our paper fills this gap through comprehensive algorithmic experiments across five VRA-covered Southern states, testing 105 configurations with four compactness metrics and three VRA metrics. We provide the first systematic Pareto frontier analysis of the VRA-compactness relationship, revealing that assumed tradeoffs are often artifacts of suboptimal algorithms rather than inherent geometric constraints. Non-MM districts generally \emph{benefit} from demographic-aware redistricting, some states achieve win-win outcomes, and geographic feasibility thresholds define the boundary of algorithmic optimization.
